\chapter{Аудит нарушителей энтропийной масштабной орбиты}
\label{ch:v10-entropy-scale-orbit-breakers}
Предыдущий мост оставил общую степень $\ell^{-4}$. Проверены десять
кандидатов, способных добавить вклад другой степени: тепловой и Casimir-
секторы ячейки, $R^2$, $\Lambda/G$, время корреляции ванны, химический
потенциал, массовая щель, dimensional transmutation, нелокальная память и
внешняя линейка.

Первые четыре внутренних кандидата снова имеют степень $-4$ и не нарушают
орбиту. Фиксированные время ванны или химический потенциал создают степень
$-3$; независимая массовая плотность создаёт степень $0$. Такие вклады
действительно могли бы выбрать длину, например
\begin{equation}\frac{A}{\ell^4}=\frac{B}{\ell^3}
\Longrightarrow \ell=\frac AB.
\end{equation}
Но $B$ тогда содержит внешний временной, энергетический или длиновой якорь,
которого текущий родитель не выводит.

Матрица десяти кандидатов по шести критериям имеет ранг шесть, однако полных
проходов нет. Лучшие кандидаты --- время корреляции ванны и химический
потенциал --- получают $5/6$, проваливая происхождение независимого масштаба.
\begin{theorem}[Граница внутреннего конформного сектора]
Все доступные внутренние локальные плотности сохраняют степень $\ell^{-4}$.
Нарушение орбиты требует дополнительного размерного свойства резервуара.
\end{theorem}
\begin{nogocriterion}[Переименование внешнего масштаба]
Введение $\tau_{bath}$, $\mu_{bath}$ или $L_{mem}$ не считается выводом,
пока эти величины не получены из того же KMS-родителя без внешней подстановки.
\end{nogocriterion}
\begin{protocol}\begin{enumerate}\item аудит: \textbf{$10/10$};\item полные проходы: \textbf{$0/10$};
\item условные нарушители: \textbf{$6/6$};\item внутренние нарушители: \textbf{$0/4$};
\item абсолютный масштаб: \textbf{$0/1$};\item следующий гейт: \textbf{происхождение времени корреляции ванны}.\end{enumerate}\end{protocol}
Машинный сертификат: \path{s2t_v10_cell_birth_four_volume_nonequilibrium_entropy_scale_orbit_breaker_candidate_audit_gate_results.json}.